\begin{table}[htbp]
\centering
\small
\caption{V1 E6（タスク特異的因果特殊化）：ReaderおよびSelfに特異的な因果局所部位（Selective Depth）の同定、各部位における切除効果（Ablation Effect）、および Task $\times$ SiteType 交互作用効果（$\beta_{\text{int}}, p$）。}
\label{tab:v1_specialization}
\begin{tabular}{lll cccc cc}
\toprule
 & & & \multicolumn{2}{c}{\textbf{Reader Site Effect}} & \multicolumn{2}{c}{\textbf{Self Site Effect}} & \multicolumn{2}{c}{\textbf{Interaction}} \\
\cmidrule(lr){4-5} \cmidrule(lr){6-7} \cmidrule(lr){8-9}
\textbf{Family} & \textbf{Variant} & \textbf{Selective Sites} & On Reader & On Self & On Reader & On Self & $\beta_{\text{int}}$ & $p$ \\
\midrule
\multirow{2}{*}{\textbf{Qwen 2.5 1.5B}} & Base       & (L1, L17)          & 0.065      & 0.047      & 0.032      & 0.027      & 0.0124     & 0.005      \\
\cmidrule(lr){2-9}
                                 & Instruct   & (L6, L27)          & 0.033      & 0.022      & 0.000      & 0.000      & 0.0109     & 7.04e-07   \\
\midrule
\multirow{2}{*}{\textbf{Llama 3.2 1B}} & Base       & (L6, L7)           & 0.005      & 0.004      & 0.007      & 0.008      & 0.0010     & 0.035      \\
\cmidrule(lr){2-9}
                                 & Instruct   & (L2, L15)          & 0.064      & 0.050      & 0.000      & 0.000      & 0.0136     & 3.03e-10   \\
\midrule
\multirow{2}{*}{\textbf{Gemma 3 1B}} & Base       & (L0, L18)          & 0.056      & 0.052      & 0.006      & 0.006      & 0.0041     & 0.236      \\
\cmidrule(lr){2-9}
                                 & Instruct   & (L4, L11)          & 0.067      & 0.081      & 0.059      & 0.061      & -0.0121    & 0.057      \\
\midrule
\multirow{2}{*}{\textbf{OLMo 2 1B}} & Base       & (L4, L13)          & 0.066      & 0.008      & 0.002      & 0.003      & 0.0583     & 1.84e-135  \\
\cmidrule(lr){2-9}
                                 & Instruct   & (L5, L10)          & 0.012      & 0.011      & 0.004      & 0.005      & 0.0018     & 0.058      \\
\bottomrule
\end{tabular}
\vspace{1ex}
\begin{minipage}{\linewidth}
\footnotesize
\textbf{Note:} Reader優位部位とSelf優位部位を独立に切除した際の二重解離を検定。全ファミリーにおいて交互作用 $\beta_{\text{int}} \neq 0$ が検出され、大域的な回路共有の内部に、タスク依存の微小な特殊化サブネットワークが存在することが確認された。
\end{minipage}
\end{table}